\documentclass[11pt]{article}
\usepackage[margin=1in]{geometry}
\usepackage{booktabs}
\usepackage[hidelinks]{hyperref}
\title{Turn clustering: a calibrated encoder comparison}
\author{Mash experimental research}
\date{6 October 2026}
\begin{document}
\maketitle
\section{Procedure}
We compare frozen all-MiniLM-L6-v2 representations with Voyage 4 and Voyage
Code 4. The matched cohort contains 156 development histories after four
credential-screening exclusions. A deterministic SHA-256 ordering assigns
104 histories to calibration and 52 to evaluation. Previously inspected
development labels are provisional agent judgments, not a pristine final test;
the original confirmation remains sealed. Each representation receives the
same 48 random-sampler Optuna trials and parameter draws, optimizing only
calibration F1. BIRCH radius, absolute membership cutoff, relative-best margin
and view are calibrated. Fitting and smoothing remain within each thread.

The local baseline uses quantized MiniLM and equal-role pooling of short
structural windows. Each remote model is tested using normalized raw floats,
the original signed-int8 codec, and complete-turn floats. Complete-turn inputs
include all current prompts and visible assistant segments with role markers.
They are split losslessly at paragraph/line boundaries if exceeding a
conservative 28,000 UTF-8-byte budget; oversize pieces have a normalized-mean
pooled view. This intervention also changes role weighting and is not a pure
context-length ablation. No LLM summaries or other-thread context are used.
Provider-default task mode, 1024 dimensions, denied data collection and disabled
fallbacks are held fixed. Raw floats and model-specific exact-text cache keys
are retained. The first remote experiment cached only int8-decoded vectors.

\section{Evaluation results}
All entries are percentages on the 52 evaluation histories. Long histories
have at least eight turns; distant pairs have turn gap at least four within
long histories. Models/configurations are not selected using these results.
\begin{center}\small
\begin{tabular}{lrrrrr}
\toprule
Representation & Precision & Recall & F1 & Long F1 & Distant recall\\
\midrule
minilm & 82.00 & 84.97 & 83.46 & 77.21 & 57.41 \\
code float & 78.03 & 90.16 & 83.65 & 77.25 & 68.52 \\
code int8 & 77.93 & 89.64 & 83.37 & 76.72 & 66.67 \\
code complete & 80.61 & 81.87 & 81.23 & 72.82 & 61.11 \\
general float & 80.37 & 91.19 & 85.44 & 80.35 & 74.07 \\
general int8 & 80.45 & 91.71 & 85.71 & 80.87 & 74.07 \\
general complete & 83.66 & 87.56 & 85.57 & 80.18 & 70.37 \\
\bottomrule
\end{tabular}
\end{center}
The companion Markdown report records fixed-view and codec ablations, selected
parameters, API charges and paired 2,000-resample thread-bootstrap intervals.
Intervals are exploratory and unadjusted for multiple comparisons. Validation
checks response ordering, dimensions, norms, ownership, source preservation and
frozen-control predictions. MLflow stores parameters, hashes and aggregate
metrics only; raw histories and detailed labels remain private.

\section{Analysis}
Calibration substantially improves the remote float pipelines relative to
using the local model's settings. Float/int8 fixed-setting differences are
small, so quantization is not supported as the main explanation for the
original regression. Chunk quality differs by model: the general-purpose
Voyage 4 chunk view is stronger than Code 4's here. A naive union can trade
precision away for recall; calibrating the BIRCH radius improves that balance.
Complete-turn input is not a universal improvement: it loses quality for
Code 4 and trades recall for precision for Voyage 4. All remote F1-gain
bootstrap intervals include zero. Voyage 4's distant-recall point gains are
promising but provisional; the frozen local baseline remains unchanged.

An off-by-one harness error counted the initially queued baseline trial as an
execution. Detected after the first local evaluation and before remote-arm
evaluation, it was corrected by restarting calibration from scratch on the
same split/cache; all final arms complete 48 trials. Original records and
source are archived. A three-input diagnostic repeat-call probe found small
Code 4 float-vector changes across differently batched calls; their cause is
not established. Float/int8 arms share the same raw vectors. No quality trial
was added by that probe. Provider repeat variability is not included in the
thread-bootstrap uncertainty estimates.

\section{Baseline selection after development evaluation}
On 6 October 2026, the user selected the calibrated Voyage 4 joint-turn float
configuration as the new experimental reference, named
\texttt{voyage4-joint-v1}. Its BIRCH radius is 0.50, absolute membership
cutoff 0.75 and relative-best margin 0.05, retaining original pooled-view
smoothing and the measured union with unsmoothed joint-piece memberships.
The source revision, matrix catalogue, annotation catalogue, report and split
hashes are pinned in the public baseline manifest. Cached matrices define
the exact reference because remote model weights cannot be pinned.

This selection uses the reported evaluation outcomes; consequently all 156
matched histories, including the former 52-thread evaluation subset, are now
development data for subsequent comparisons. The original confirmation set
remains sealed. Selection freezes a useful experimental reference rather than
claiming a statistically established superiority or production acceptance.
The historical MiniLM reference is preserved and local production search
continues using its existing model. No new inference or confirmation evaluation
was performed as part of the freeze.

\section{Sources and limitations}
The frozen study remains a reference and production behavior is unchanged.
Remote model weights cannot be pinned, and repeated API calls may expose
provider drift. Calibration/holdout division cannot undo earlier inspection
of this development corpus. Human-reviewed, independently held-out evaluation
is still required before strong quality claims.
\begin{thebibliography}{9}
\bibitem{minilm} Sentence Transformers. all-MiniLM-L6-v2 model card.
\url{https://huggingface.co/sentence-transformers/all-MiniLM-L6-v2}
\bibitem{voyage} Voyage AI. Text embeddings: models and API task modes.
\url{https://docs.voyageai.com/docs/embeddings}
\bibitem{code} Voyage AI. Voyage-code-4: code retrieval built for coding agents.
\url{https://blog.voyageai.com/2026/08/13/voyage-code-4/}
\bibitem{optuna} T. Akiba et al. Optuna: A Next-generation Hyperparameter
Optimization Framework. KDD, 2019. \url{https://arxiv.org/abs/1907.10902}
\bibitem{mlflow} MLflow. Tracking APIs.
\url{https://mlflow.org/docs/latest/ml/tracking/tracking-api/}
\end{thebibliography}
\end{document}
